% !TEX root = ../../../main.tex

\subsection{数据资源与技术选择}

实时单词识别模型主要使用系统自采样本进行训练。
管理员在系统维护过程中可通过样本采集功能为指定手势动作采集训练样本，Python 手势识别服务接收采集帧并保存为 raw dataset。
选择自采数据作为实时单词识别的数据来源，是因为该任务需要直接服务于 HearBridge 当前维护的手语资源和训练动作，自采样本更贴近系统实际训练项、采集方式和运行环境。

在技术选择上，实时单词识别采用基于关键点特征的固定词表分类方案。
识别服务使用 MediaPipe 对手部和人体姿态关键点进行提取，将原始图像帧转换为结构化时序数据，再通过特征转换流程生成模型训练所需的数据。
模型训练部分采用 TensorFlow/Keras 实现固定词表分类模型，具体使用 1D CNN 对固定长度的时序特征进行分类。
